\section{Training: Value-Based RL as the Spine}
\label{sec:training}

\subsection{Off-policy Q-learning core}
\label{sec:training-core}

The core is approximate Q-learning over sampled transitions
$(s, a, r, s')$ with $a = (\textit{verb}, \textit{target})$:
\[
y \;=\; r \;+\; \gamma \max_{(\textit{v}', \textit{t}') \in A(s')}
Q_{\mathrm{target}}(s', \textit{v}', \textit{t}'),
\qquad
\mathcal{L}_Q = \big(Q(s, \textit{verb}, \textit{target}) - y\big)^2,
\]
with a replay buffer and a periodically synced target network. The roles
are: the verb and pointer heads \emph{act} (they are the deployed policy),
and $Q$ \emph{criticizes}: a masked softmax over $Q$ at temperature $\tau$
supplies the distribution toward which the heads are trained. A learned
$Q$ is the natural object in this setting: it folds the expensive
expectation over LLM transitions into a cheap action value, so at
deployment the controller selects an action without sampling any LLM
transition to look ahead. Off-policy learning is load-bearing: the same
buffer accepts random exploration, teacher trajectories
(\cref{sec:training-ladder}), and search-generated experience
(\cref{sec:training-mcts}) with no change to the update. Under
terminal-only reward, one-step bootstrapping is fragile; we regress on
n-step and Monte Carlo returns from the verifier terminal, which requires
no search in the loop. The price of doing so without correction must be
stated rather than hidden: uncorrected n-step and Monte Carlo returns
estimate $Q^{\mu}$, the value of the behavior policy $\mu$ that generated
the trajectory (the teacher, an exploring variant of the current heads, or
search), not $Q^{*}$. We accept this bias in v1, on the grounds that the
behavior mixture tracks the improving policy closely and that under sparse
terminal reward long uncorrected returns carry the only usable signal;
importance-weighted or truncated corrections (Retrace-style) are the
sanctioned upgrade if the bias proves material.

\subsection{The three-rung ladder}
\label{sec:training-ladder}

Learning a policy over a variable, growing action space under sparse reward
is the highest-risk element of the design; a training ladder removes the
cold start. Two LLM roles are distinct throughout: the \emph{generator} (the
small frozen LLM that produces content under the lens; the deployed model)
and the \emph{controller-teacher} (an LLM, possibly much larger, used only
at training time to pick $(\textit{verb}, \textit{target})$ actions).

\begin{enumerate}[leftmargin=1.5em, itemsep=1pt]
\item \textbf{LLM-as-controller.} The teacher picks actions; the small
frozen generator fills content under the lens; the verifier checks. This
validates the full interface before any training and produces bootstrap
trajectories. The teacher orchestrates the \emph{small} generator during
collection, so trajectories reflect the content the deployed generator can
actually produce: decision quality from the teacher, state distribution from
the generator.
\item \textbf{Behavior-cloning warm start.} The verb and pointer heads are
trained to imitate the teacher's actions over the global state embedding.
\item \textbf{RL fine-tune.} The off-policy core of
\cref{sec:training-core} runs from the warm-started heads, consuming both
teacher and self-generated experience.
\end{enumerate}

This is distillation of a different object than in o1/R1-style pipelines:
the big model's weights are never touched and its chain of thought is never
the product; what is distilled is its \emph{policy over typed operations},
into a controller small enough to be negligible at deployment.

\subsection{MCTS as a strictly additive option}
\label{sec:training-mcts}

Because the core is off-policy, Monte Carlo Tree Search contributes in two
purely additive ways: as a better behavior policy (search-guided
exploration adds higher-quality rows to the same buffer) and as a source of
better targets (chance-node estimates $\hat{Q}_{\mathrm{MCTS}}(s,a)$ replace
one-step bootstraps; visit-count distributions replace softmax-over-Q as
policy targets). Distillation from visit counts, the spine of classical
expert iteration, is here exactly this add-on: the MCTS-on variant of the
same losses, evaluated as an ablation. At deployment, a budget
$B \in \{0, 1, 2, \dots\}$ selects between acting from the learned heads
directly ($B{=}0$) and running $B$ PUCT simulations from the current state
with the learned $Q$ and policy as priors, which is the system's test-time
compute knob. \stub{final hyperparameters: buffer size, n-step length,
target-sync period, $\tau$ schedule, $c_{\mathrm{puct}}$}
